\scanpage{47}
racines de $\Sigma$, sauf une qui est transformée en $-\psi$ ($\psi=$ plus grande racine).

Déterminer aussi le normalisateur de $\unsure{\kappa|K}$, et l'action de $\mathrm{Norm}(\kappa)/\mathrm{Zent}(\kappa)$ sur $\mu^{\unsure{\circ}}_h$ ; trouve-t-on tout $\mathrm{Aut}(\mu_h)=(\mathbb Z/h\mathbb Z)^{*}$ ?

\margin{$G$ simple ?}
(iii) Soit $\deleted{\unsure{\kappa}}$
\[
\kappa:\mu_h\unsure{-|\to}T
\]
un hom principal. Considérons comme dessus ${\mathfrak g}^{1}\inserted{=\sum_{\alpha\in Q}\mathfrak g^{\alpha}\ \text{(qui détermine $\kappa$)}}$. Kostant prouve (en car. $0$)

a) Toute section $\sum x^{\alpha}\in\Gamma\mathfrak g^1$ $(\unsure{x^{2}|x^{\alpha}}\in\Gamma\mathfrak g^{\alpha}]$ telle que les $x^{\alpha}$ ne s'annulent pas est \uline{$1/2$ simple régulière}.
% slanted
\margin{pourrait marcher au mieux si $h$ est inv. sur $S$}

b) Si un des $x^{\alpha}$ est nulle, elle est \uline{nilpotente}.
% slanted
\margin{marche sans doute $\overline{\mathrm{car}}$. Elle est nilpotente régulière ssi tous les autres ne s'annulent pas ?}

c) Soient deux droites $D,D'\subset\check V(\mathfrak g^{1})^{\mathrm{ss}}$, qui soient du cas a). Alors $\exists$\inserted{loc. fpqc} $t\in\Gamma T$ telle que $D'=\mathrm{ad}(t).D$

[Sans doute le stabilisateur dans $T$ d'un $D$ comme dessus est $\kappa(\mu_h)$, donc $t$ est une section canonique de $T/\kappa(\mu_h)$. On peut dire aussi que \deleted{\unsure{tout forme les isom.}}\inserted{dans} $\underline{\mathrm{Isom}}_{\underline{0}}(D\overset{\sim}{\to}D')$ \deleted{\unsure{qui envoient} \illegible{1}, \unsure{classe de}\illegible{1}}\inserted{on peut définir, un sous-}\deleted{\unsure{fibré}} $\mu_h$-torseur canonique [savoir \deleted{\unsure{l'ens des}} le fibré des $t\in T$ qui envoient $D$ dans $D'$], et pour un isom donné de $D$ avec $D'$ dans cette classe, il $\exists 1$ (unique) $t\in T$ qui l'induit [Cette information résulte du fait que un élément de $T$ qui est l'identité sur $D$ est la section $1$].
% archive sheet number: 46

